\documentclass{article}
\usepackage[spanish]{babel}
\usepackage[utf8]{inputenc}
\usepackage{amsmath,amssymb}
\usepackage{graphicx}
\usepackage{hyperref}

\title{Análisis de Complejidad del Problema del Comerciante Holandés}
\author{}
\date{}

\begin{document}

\maketitle

\begin{abstract}
Este documento presenta un análisis de complejidad computacional del Problema del Comerciante Holandés, demostrando su NP-dureza mediante una reducción desde el problema Prize-Collecting TSP.
\end{abstract}

\section{Introducción}
\label{sec:introduccion}

El Problema del Comerciante Holandés (CH) es un problema de optimización combinatoria que modela las decisiones estratégicas de ruta y transacción de un comerciante marítimo con múltiples restricciones. En este documento demostramos que el problema de decisión asociado es \textbf{NP-Completo}. Esta demostración se estructura en dos partes: (1) estableciendo su \textbf{NP-dureza} mediante una reducción polinomial desde el problema Prize-Collecting TSP (PCTSP), y (2) probando su \textbf{pertenencia a la clase NP} mediante un verificador polinomial. Este resultado fundamental implica que no es posible resolver el problema de forma exacta y eficiente (en tiempo polinomial) para instancias de tamaño arbitrario, a menos que P = NP, justificando así el desarrollo de algoritmos aproximados y heurísticos en fases posteriores del proyecto.

\section{Demostración de NP-Dureza} 
\label{sec:np-hardness}

\subsection{Definición Formal del Problema del Comerciante Holandés (CH)}
\label{subsec:def_ch}

Una instancia del problema \textsc{ComercianteHolandés} está definida por la tupla $(G, r, \mathcal{M}, \mathcal{P}, B, T_{max}, f)$ donde:

\begin{itemize} 
    \item $G=(V, E)$ es un grafo completo no dirigido que representa la red de puertos. $V = \{v_0, v_1, ..., v_n\}$, con $v_0$ siendo el puerto de origen (Ámsterdam). 
    \item $r \in \mathbb{R}^+$ es el capital inicial.
    \item $\mathcal{M}$ es el conjunto de $k$ tipos de mercancías. 
    \item $\mathcal{P} = \{ (c^{compra}_{v,m}, c^{venta}_{v,m}) \mid v \in V \setminus \{v_0\}, m \in \mathcal{M} \}$ define el precio de compra y venta de cada mercancía en cada puerto. 
    \item $B \in \mathbb{N}$ es la capacidad máxima de carga del barco (en unidades de mercancía). 
    \item $T_{max} \in \mathbb{R}^+$ es el tiempo máximo total permitido para la expedición. 
    \item $f: E \rightarrow \mathbb{R}^+$ asigna a cada arista $e=(u,v)$ un tiempo de viaje $t_{uv}$ y un costo operativo $c_{uv}$. 
\end{itemize}

Una \textbf{solución factible} es un ciclo simple $(v_0, v_{i_1}, ..., v_{i_l}, v_0)$ y una secuencia de decisiones de compra/venta en cada puerto visitado, tal que: 
1) No se supera la capacidad $B$ en ningún momento. 
2) El capital nunca es negativo tras cubrir costos operativos. 
3) El tiempo total (viaje + operaciones) $\leq T_{max}$.

El \textbf{objetivo} es maximizar el capital final al regresar a $v_0$.

\subsection{Definición del Problema Reducido: \textsc{CH-Restringido}}
\label{subsec:def_ch_rest}

Para la reducción, definimos una versión fuertemente restringida de \textsc{ComercianteHolandés}, denominada \textsc{CH-Restringido}. Esta instancia se obtiene fijando o eliminando los siguientes parámetros y decisiones del problema original:

\begin{enumerate} 
    \item \textbf{Mercancías}: Se fija $|\mathcal{M}| = 1$. Solo existe una mercancía genérica. 
    \item \textbf{Precios y Transacciones}: Se fijan $c^{compra}_{v} = 0$ y $c^{venta}_{v} = p_v \geq 0$ para todo puerto $v \in V \setminus \{v_0\}$. Esto implica que al visitar $v$, se obtiene automáticamente una ganancia neta fija $p_v$, sin que medie decisión de compra ni consumo de capacidad. 
    \item \textbf{Capacidad}: Se fija $B = \infty$. La capacidad deja de ser una restricción. 
    \item \textbf{Restricción Financiera}: Se elimina la condición de capital no negativo intermedio. El capital es simplemente $r + \sum p_v - \sum c_{uv}$. 
    \item \textbf{Tiempo}: Se fija $T_{max} = \infty$ o a un valor suficientemente grande que no restrinja ningún tour posible. Los costos de arista $c_{uv}$ representan \textit{penalizaciones} (pérdida de capital) por el viaje. 
\end{enumerate}

Por lo tanto, una instancia de \textsc{CH-Restringido} se define simplemente por $(G, v_0, \{p_v\}_{v \in V}, \{c_{e}\}_{e \in E})$, y su objetivo es encontrar un ciclo que comience y termine en $v_0$, visite un subconjunto $S \subseteq V \setminus \{v_0\}$, y maximice $\sum_{v \in S} p_v - \sum_{e \in \text{tour}} c_{e}$.

\subsection{Reducción Polinomial desde \textsc{Prize-Collecting TSP}}
\label{subsec:reduccion_pctsp}

\subsubsection*{Definición de PCTSP} 
Una instancia del problema \textsc{PrizeCollectingTSP} (PCTSP) está dada por $(H, s, \{\pi_u\}_{u \in U}, \{\omega_{ij}\}_{(i,j) \in F})$, donde $H=(U, F)$ es un grafo completo, $s \in U$ es el nodo raíz, $\pi_u \geq 0$ es el premio por visitar el nodo $u$, y $\omega_{ij} \geq 0$ es el costo de traversar la arista $(i,j)$. El objetivo es encontrar un ciclo que comience y termine en $s$, visite un subconjunto de nodos, y maximice la suma de premios menos la suma de costos del tour.

\textbf{Teorema 1 (NP-Dureza de PCTSP)}: \textsc{PrizeCollectingTSP} es NP-duro. Esto se deduce directamente de ser una generalización del clásico \textsc{TravelingSalesmanProblem} (TSP, NP-duro \cite{karp1972}) donde los premios de todos los nodos excepto $s$ son mayores que la suma máxima posible de costos.Además la NP-dureza de este problema está bien establecida en la literatura, tal como se documenta en estudios sobre problemas de rutas con recolección de premios \cite{MarchettiSpaccamela2007PrizeCollectingTS}.
\cite{MarchettiSpaccamela2007PrizeCollectingTS}.
\subsubsection*{Construcción de la Reducción} 
Dada una instancia arbitraria $I_{P} = (H, s, \{\pi_u\}, \{\omega_{ij}\})$ de PCTSP, construimos una instancia $I_{C} = (G, v_0, \{p_v\}, \{c_{e}\})$ de \textsc{CH-Restringido} en tiempo polinomial de la siguiente manera: 
\begin{enumerate} 
    \item Para cada nodo $u \in U$ en $H$, creamos un puerto correspondiente $v \in V$ en $G$. El nodo raíz $s$ se mapea al puerto de origen $v_0$. 
    \item El grafo $G$ es completo. Para cada arista $(i,j)$ en $H$ con costo $\omega_{ij}$, definimos el costo de viaje/costo operativo en $G$ como $c_{v_i v_j} = \omega_{ij}$. 
    \item Para cada puerto $v$ correspondiente a un nodo $u \neq s$, definimos la ganancia fija $p_v = \pi_u$. 
\end{enumerate} 
Esta construcción se realiza en $O(|U|^2)$.

\subsubsection*{Equivalencia de Soluciones y Optimalidad} 
Existe una correspondencia biyectiva entre los tours factibles de $I_P$ y $I_C$ que preserva la función objetivo. Dado un tour $T_P = (s, u_{i_1}, ..., u_{i_m}, s)$ para $I_P$ con valor $\sum \pi_{u_{i_k}} - \sum \omega_{tour}$, el tour correspondiente $T_C = (v_0, v_{i_1}, ..., v_{i_m}, v_0)$ en $I_C$ tiene un capital final de $r + \sum p_{v_{i_k}} - \sum c_{tour}$. Dado que $r$ es constante, maximizar el capital final en $I_C$ es equivalente a maximizar $\sum p_{v_{i_k}} - \sum c_{tour}$, que es idéntico a la función objetivo de $I_P$. El argumento inverso es análogo.

\textbf{Lema 1}: Existe un tour en $I_P$ con valor $\geq K$ si y solo si existe un tour en $I_C$ con capital final $\geq r + K$.

\subsubsection*{Conclusión de la Reducción} 
Dado que PCTSP es NP-duro (Teorema 1) y hemos construido una reducción en tiempo polinomial que preserva la optimalidad desde PCTSP a \textsc{CH-Restringido}, concluimos que:

\textbf{Teorema 2}: \textsc{ComercianteHolandés-Restringido} es NP-duro.

\subsection{Preservación de la Complejidad al Reintroducir Restricciones}
\label{subsec:preservacion_complejidad}

El Teorema 2 demuestra la NP-dureza de una versión simplificada ($\Pi_{restr}$) de nuestro problema original ($\Pi_{orig}$). Para concluir definitivamente que $\Pi_{orig}$ es NP-duro, debemos analizar que reintroducir cada una de las restricciones eliminadas no transforma el problema en uno más fácil (en P). Demostraremos que, para cada restricción, $\Pi_{restr}$ es un \textbf{caso particular} de $\Pi_{orig}$ bajo una configuración específica, por lo que $\Pi_{orig}$ es al menos tan difícil.

Sea $\Pi_{orig}(X, Y, Z, ...)$ el problema general con restricciones $X, Y, Z, ...$, y $\Pi_{restr}$ el problema sin ellas. Si $\Pi_{restr}$ es NP-duro, y existe una asignación de valores para $X, Y, Z, ...$ tal que $\Pi_{orig}(val_X, val_Y, val_Z, ...) \equiv \Pi_{restr}$, entonces $\Pi_{orig}$ es NP-duro.

\begin{enumerate}
    \item \textbf{Múltiples Mercancías ($|\mathcal{M}| \geq 1$)}. En $\Pi_{restr}$, $|\mathcal{M}|=1$. Sea una instancia de $\Pi_{orig}$ con $|\mathcal{M}|=k>1$ donde: i) los precios de compra de las mercancías $2, ..., k$ son superiores a cualquier ganancia posible, haciéndolas no rentables; ii) los precios de venta de la mercancía 1 son idénticos a $p_v$; y iii) la capacidad es suficiente para cargar solo la mercancía 1. Cualquier solución óptima ignorará las mercancías $2, ..., k$ y se comportará como en $\Pi_{restr}$. Por tanto, $\Pi_{restr}$ es un caso particular de $\Pi_{orig}$ con múltiples mercancías, que no reduce la complejidad.

    \item \textbf{Capacidad Finita de Carga ($B \in \mathbb{N}$)}. En $\Pi_{restr}$, $B = \infty$. Fijar $B$ a un valor suficientemente grande (e.g., $B \geq \sum_{v \in V} 1$) en $\Pi_{orig}$ replica las condiciones de capacidad ilimitada, ya que nunca se alcanzará el límite con una única mercancía de ganancia unitaria. Reducir $B$ restringe el espacio de soluciones, añadiendo una dimensión de dificultad equivalente a un problema de la mochila (Knapsack), el cual es NP-duro. La adición de una restricción NP-dura no puede simplificar un problema que ya es NP-duro.

    \item \textbf{Restricción Financiera Intermedia (Capital no Negativo)}. En $\Pi_{restr}$ esta restricción se omite. En $\Pi_{orig}$, si el capital inicial $r$ es mayor que la suma máxima posible de costos de viaje, la restricción nunca se activa, replicando $\Pi_{restr}$. Reducir $r$ introduce un estado financiero dinámico que debe ser verificado en cada puerto, aumentando la complejidad del espacio de estados, no reduciéndola.

    \item \textbf{Límite Temporal Estricto ($T_{max} \in \mathbb{R}^+$)}. En $\Pi_{restr}$, $T_{max} = \infty$. Configurar $T_{max}$ mayor que la duración del tour más largo posible en $\Pi_{orig}$ hace que la restricción sea inactiva, replicando $\Pi_{restr}$. Imponer un límite temporal estricto convierte el problema en una variante del "Orienteering Problem" (también NP-duro), añadiendo una capa de dificultad adicional.

    \item \textbf{Decisiones de Compra/Venta Complejas}. En $\Pi_{restr}$, la ganancia es automática ($c^{compra}_v = 0, c^{venta}_v = p_v$). En $\Pi_{orig}$, fijar estos precios de forma idéntica y definir que el beneficio se obtiene al llegar al puerto (sin consumo de capacidad) simula el comportamiento de $\Pi_{restr}$. Permitir precios variables y decisiones estratégicas introduce un problema de arbitraje y gestión de inventario, que es una generalización estricta y más difícil.
\end{enumerate}

\textbf{Conclusión del Análisis}: Para cada restricción $R_i$, existe una configuración de $\Pi_{orig}$ que la hace equivalente a $\Pi_{restr}$. Cualquier ajuste que active genuinamente $R_i$ (e.g., reducir $B$, reducir $T_{max}$) \textbf{restringe} el espacio de soluciones factibles y/o añade una dimensión de optimización NP-dura. Por lo tanto, $\Pi_{orig}$ no es más simple que $\Pi_{restr}$; es una generalización que preserva y potencialmente incrementa la complejidad.

\textbf{Teorema 3 (NP-Dureza del Problema General)}: El problema \textsc{ComercianteHolandés} ($\Pi_{orig}$) es NP-duro.
\textbf{Demostración}: Directa por el Teorema 2 y el análisis anterior, que establece que $\Pi_{restr}$ es un caso particular de $\Pi_{orig}$. \hfill $\square$

\subsection{Demonstración de Pertenencia a la Clase NP} \label{subsec:
pertenencia_np}

Para establecer que el \textsc{ComercianteHolandés} es NP-Completo, además de
ser NP-duro (Teorema 3), debemos demostrar que pertenece a la clase NP. Es decir,
que dada una supuesta solución al problema, podemos verificar su factibilidad y
calcular su valor objetivo (el capital final) en tiempo \textbf{polinomial}
respecto al tamaño de la entrada.

\subsubsection*{Certificado y Verificador} Un certificado $C$ para una
instancia de $\Pi_{orig}$ debe incluir: \begin{enumerate} \item Una secuencia de
puertos (un tour) $T = (v_0, v_{i_1}, v_{i_2}, ..., v_{i_m}, v_0)$. \item Para
cada puerto visitado $v_{i_j}$, una lista de operaciones de compra y venta por
tipo de mercancía. \end{enumerate}

El tamaño de este certificado está acotado polinomialmente: el tour tiene como
máximo $|V|$ nodos, y las operaciones en cada puerto son para $|\mathcal{M}|$
mercancías.

Un algoritmo verificador $V(I, C)$ debe comprobar los siguientes puntos en
tiempo polinomial: \begin{enumerate} \item \textbf{Validez del Tour}: Verificar
que $T$ es un ciclo que comienza y termina en $v_0$, y que ningún puerto interno
se repite. Esto se hace en $O(|V|)$. \item \textbf{Cumplimiento de la Capacidad}:
Para cada segmento entre puertos, calcular la carga total a bordo sumando las
compras y restando las ventas. Verificar que en ningún punto excede $B$. Esto
requiere recorrer la secuencia de operaciones una vez: $O(m \cdot
|\mathcal{M}|)$. \item \textbf{Cumplimiento de la Restricción Financiera}:
Simular el flujo de capital. Partiendo de $r$, para cada puerto: restar costos
de viaje, sumar ganancias por ventas, restar costos por compras, y verificar que
el capital no sea negativo. También es $O(m \cdot |\mathcal{M}|)$. \item
\textbf{Cumplimiento del Límite Temporal}: Sumar los tiempos de viaje $t_{uv}$
para todas las aristas en $T$ y los tiempos de operación en puertos. Verificar
que el total no exceda $T_{max}$. $O(m)$. \item \textbf{Cálculo del Objetivo}:
El valor del certificado es el capital final tras la última operación. Se
obtiene como salida del paso 3. \end{enumerate}

Cada una de estas operaciones es una suma, resta o comparación que se ejecuta
una cantidad de veces proporcional al tamaño del certificado. Por lo tanto, el
tiempo total de verificación es polinomial en el tamaño de la entrada $I$ (que
incluye $|V|$, $|E|$, $|\mathcal{M}|$).

\textbf{Lema 3 (Pertenencia a NP)}: Existe un verificador $V(I, C)$ que decide
en tiempo polinomial si un certificado $C$ es una solución factible para
$\Pi_{orig}$ y cuál es su valor. Por lo tanto, $\Pi_{orig} \in \text{NP}$.

\subsubsection*{Conclusión de NP-Completitud} Combinando el resultado de
NP-dureza (Teorema 3) con la pertenencia a NP (Lema 3), podemos establecer la
clasificación final del problema.

\textbf{Teorema 4 (NP-Completitud)}: El problema de decisión asociado al
\textsc{ComercianteHolandés} es NP-Completo. \textbf{Demostración}: El problema
de decisión pregunta: "¿Existe un plan de ruta y transacciones con capital final
$\geq K$?". Por el Lema 3, este problema está en NP. Por el Teorema 3, es
NP-duro. En consecuencia, es NP-Completo. \hfill $\square$

Esta clasificación justifica formalmente el empleo de técnicas de diseño
algorítmico no exactas (heurísticas, aproximaciones) en las siguientes fases del
proyecto, como se solicita en la Fase 3.

\section{Conclusiones}
\label{sec:conclusiones}

En este trabajo se ha realizado un análisis formal de la complejidad computacional del Problema del Comerciante Holandés. Los resultados principales son:
\begin{itemize}
    \item Se definió una versión restringida del problema (\textsc{CH-Restringido}) y se demostró que es NP-dura mediante una reducción polinomial desde el problema Prize-Collecting TSP (Teorema 2).
    \item Se argumentó de manera rigurosa, restricción por restricción, que la complejidad NP-dura se preserva al reintroducir todas las características del problema general (Teorema 3).
    \item Se demostró que el problema de decisión asociado pertenece a la clase NP, presentando un certificado y un verificador polinomial (Lema 3 y Teorema 4).
\end{itemize}
En consecuencia, se concluye que el \textbf{problema de decisión del Comerciante Holandés es NP-Completo}. Este hallazgo justifica teóricamente el enfoque práctico del proyecto, que en las siguientes fases se centrará en el diseño, implementación y evaluación empírica de algoritmos de aproximación y heurísticas para abordar instancias realistas del problema.

\bibliographystyle{plain}
\bibliography{referencias}

\end{document}